\subsection{TOPOLOGICAL GROUPS WITH OPERATORS}

On a set G, a structure of a group with operators and a topology are said to be compatible if the topology and the group structure of G are compatible (§ 1, no. 1) and if in addition the endomorphisms of G produced by the operators are continuous. The set G, together with the given topology and structure of group with operators, is then said to be a topological group with operators.

If H is a stable subgroup of a topological group G with operators, then the topology induced on H by the topology of G is compatible with the structure of group with operators on H. Furthermore:

PROPOSITION I. If H is a stable subgroup of a topological group G with operators, then the closure \(\overline{H}\) of H in G is a stable subgroup of G.

We know already (§ 2, no. 1, Proposition 1) that \(\overline{\mathbf{H}}\) is a subgroup of G. Also if \(\alpha\) is any operator on G, the image of H under the continuous mapping \(x \to x^{\alpha}\) is contained in H, and therefore the image of \(\overline{\mathbf{H}}\) is contained in \(\overline{\mathbf{H}}\) (Chapter I, § 2, no. 1, Theorem 1).

Let H be a stable normal subgroup of a topological group G with operators. Then for each operator \(\alpha\) on G, the mapping of G/H into itself induced by \(x \rightarrow x^{\alpha}\) is continuous (§ 2, no. 8, Remark 3), and the structure of group with operators on G/H is therefore compatible with the quotient topology on G/H.

Let \((\mathbf{G}_i)_{i\in I}\) be a family of topological groups with operators, where each \(\mathbf{G}_i\) is assumed to have the same operator set \(\Omega\) . For each \(\alpha \in \Omega\) , the mapping \(x\to ((\mathrm{pr}_i x)^{\alpha})\) of \(\mathbf{G} = \prod_{i\in I}\mathbf{G}_i\) into itself is continuous (Chapter I, § 4, no. 1, Proposition 1), and the structure of groups with operators on \(\mathbf{G}\) is therefore compatible with the product topology on \(\mathbf{G}\) .

If G is a Hausdorff topological group with operators and if G has a completion \(\hat{G}\) (§ 3, no. 4), then every endomorphism \(x \to x^{\alpha}\) of G defined by an operator on G can be extended by continuity to an endomorphism of \(\hat{\mathbf{G}}\) (§ 3, no. 3, Proposition 5). Hence \(\hat{\mathbf{G}}\) has the structure of a topological group with operators, and the operator set is the same for \(\hat{\mathbf{G}}\) as for \(\mathbf{G}\) .

